% mode: ini-etex
% Aimed at §1071-§1076 and §1610: how the second pass finds the word to hyphenate, before hyphenate itself.
% §1073: kerns, math L nodes and whatsits are skipped, \setlanguage takes effect, a box or a math shift stops it; empty
% ligature; a capital with \uchyph off and on; lc_code 0.  §1074: another font, a non-letter, 63 and 64 letters; §1075:
% ligatures in the word, a ligature with lc_code 0, the bchar bit; font kerns in the word.  §1076: what follows the word
% decides (kern, math, box, rule, mark, penalty, insert, vadjust, whatsit).  \hyphenchar out of range; left+right above 63.
% \savinghyphcodes makes hyph_index nonzero, set_lc_code's other branch.
\input prelude
\hyphenchar\tenrm=`\-  \lefthyphenmin=2 \righthyphenmin=2
\savinghyphcodes=1
\language=0
\patterns{a1b b1c c1d d1e e1f f1g g1h h1i i1j j1k k1l l1m m1n n1o o1p p1q q1r r1s s1t t1u u1v v1w w1x x1y y1z}
\language=1 \lccode`A=`a
\patterns{1a1 1b1 1c1 1d1 1e1 1f1 1g1 1h1 1i1 1l1 1o1 1u1}
\lccode`A=`A
\language=0 \savinghyphcodes=0
\hyphenation{}
\font\tenbf=cmbx10 \hyphenchar\tenbf=`\-
\tracingparagraphs=1 \hsize=20pt \parindent=0pt \pretolerance=-1 \tolerance=10000 \hbadness=10000 \hfuzz=16383pt
\showboxbreadth=1000 \showboxdepth=3
\def\T#1{\setbox1\vbox{\noindent x #1\par}\showbox1 }
% plain word; kern, whatsit and L node before it
\T{abcdefgh}
\T{\kern0pt abcdefgh}
\T{\special{}abcdefgh}
\TeXXeTstate=1 \T{\beginL abcdefgh\endL}
\TeXXeTstate=0
% an explicit kern is not a font kern, §1073 skips only subtype normal
\T{AVabcdefgh}
\T{\/abcdefgh}
% \setlanguage on the way to the word
\T{\setlanguage1 abcdefgh}
\T{\language=1 abcdefgh \language=0 abcdefgh}
% a box or a math shift stops the search
\T{\hbox{}abcdefgh}
\T{\vrule abcdefgh}
\T{$\relax$abcdefgh}
\T{\penalty0 abcdefgh}
\T{\mark{}abcdefgh}
\T{\vadjust{}abcdefgh}
% capitals
\T{Abcdefgh} \uchyph=1 \T{Abcdefgh} \uchyph=0
\T{ABCDEFGH}
% lc_code 0 skips on, then a letter
\T{1(abcdefgh}
\T{--abcdefgh}
% a ligature first, as the first thing, then inside
\T{fiabcdefgh}
\T{abcdfiefgh}
\T{abcfflefgh}
\T{abc---efgh}
\T{a{\tenbf bcdefgh}}
% §1074: another font, a digit, end of the letters
\T{abc{\tenbf def}gh}
\T{abc1efgh}
% 63 and 64 letters: hn=63 stops §1074
\T{abcdefghijklmabcdefghijklmabcdefghijklmabcdefghijklmabcdefghijk}
\T{abcdefghijklmabcdefghijklmabcdefghijklmabcdefghijklmabcdefghijkl}
\T{abcdefghijklmabcdefghijklmabcdefghijklmabcdefghijklmabcdefghijklmabc}
% a ligature at the limit
\T{abcdefghijklmabcdefghijklmabcdefghijklmabcdefghijklmabcdefghijfi}
\T{abcdefghijklmabcdefghijklmabcdefghijklmabcdefghijklmabcdefghijklfi}
% font kern inside the word, in cmr10 between some letters
\T{oxvoxvawayabcdefgh}
% §1076: what follows
\T{abcdefgh\kern0pt x}
\T{abcdefgh\/}
\T{abcdefgh\hbox{}}
\T{abcdefgh\vrule}
\T{abcdefgh$\relax$}
\T{abcdefgh\penalty0 x}
\T{abcdefgh\mark{}x}
\T{abcdefgh\vadjust{}x}
\T{abcdefgh\insert100{}x}
\T{abcdefgh\special{}x}
\T{abcdefgh\ x}
\T{abcdefgh\discretionary{}{}{}x}
\TeXXeTstate=1 \T{abcdefgh\beginL x\endL}
\TeXXeTstate=0
% too short for the minima
\lefthyphenmin=4 \righthyphenmin=4 \T{abcdefg} \T{abcdefgh}
\lefthyphenmin=40 \righthyphenmin=40 \T{abcdefgh}
\lefthyphenmin=31 \righthyphenmin=32 \T{abcdefgh}
\lefthyphenmin=100 \righthyphenmin=-3 \T{abcdefgh}
\lefthyphenmin=2 \righthyphenmin=2
% \hyphenchar out of range
\hyphenchar\tenrm=-1 \T{abcdefgh}
\hyphenchar\tenrm=256 \T{abcdefgh}
\hyphenchar\tenrm=255 \T{abcdefgh}
\hyphenchar\tenrm=`\-
% the language given by \lccode saved with the patterns: A is not a letter in language 0 here, and is in 1
\lccode`A=0
\T{\language=1 abcAdef}
\T{\language=1 Abcdef}
\lccode`A=`a
\T{\language=2 abcdefgh}
\T{\setlanguage255 abcdefgh}
% a font kern on the way to the word, between two characters that are no letters: language 2 reads \lccode
\language=2 \lccode`A=0 \lccode`V=0
\T{AVabcdefgh}
\T{Vabcdefgh}
\lccode`A=`a \lccode`V=`v \language=0
\end
